\section{Agent A: the BDI agent}
\label{sec:bdi}

\paragraph{Beliefs.} The map is treated as static known knowledge: the graph is
built once from the \texttt{map} event, with walls removed, delivery and
spawner tiles tagged, and one-way arrow tiles encoded as directed edges (an
edge is omitted when it opposes a tile's arrow). Parcels and other agents come
from \emph{partial} sensing and are stamped with a \texttt{lastSeen} time. A
parcel's reward is projected forward by subtracting whole decay ticks elapsed
since it was last seen (matching the server's whole-tick semantics, not a
fractional estimate). Negative evidence is applied: a remembered parcel that is
absent from a now-visible tile is deleted.

\paragraph{Belief reconciliation --- a live finding.} Early live testing
exposed a concrete ``my knowledge was wrong, not the world'' case. The running
server's pickup acks carried no \texttt{id} field (unlike the SDK types), so the
agent marked a phantom parcel as carried and then issued 336 futile
\texttt{putdown} retries in one 60\,s run, scoring 173. The fix normalizes all
ack ids through one layer and reconciles beliefs against the server
(\texttt{markTilePickedUp}/\texttt{clearCarried}) when no id is usable; the same
scenario then scored 800 with 30 deliveries and zero futile retries. The empty
ack shape was later confirmed on the course server too.

\paragraph{Options and strategies.} The option generator emits
\texttt{go\_pick\_up}, \texttt{deliver\_carried}, \texttt{explore} and
\texttt{wait} (plus \texttt{go\_to\_mission\_target} for Agent~B). A strategy
assigns each option a utility and \texttt{selectOption} takes the argmax, so
strategies decide but never act. Four baselines are implemented:
\texttt{greedy-nearest} (nearest parcel; deliver only when nothing is
reachable), \texttt{reward-distance} (maximize projected \emph{delivered} value,
reward minus decay\,$\times$\,distance), \texttt{delivery-threshold} (batch
until $N$ parcels, then deliver), and \texttt{mission-aware}
(\texttt{reward-distance} plus mission constraints, Agent~B's default). A
post-hoc \texttt{reward-distance-total} variant is discussed in
Section~\ref{sec:experiments}. Adding a strategy is a single file plus a
registry entry; the rest of the pipeline is untouched.

\paragraph{Intention revision.} Revision follows a replace policy with an
additive hysteresis margin: a candidate option must exceed the current
intention's utility by the margin before the agent switches, damping
target thrashing when several parcels have similar value. An intention stops
when achieved, impossible, or no longer worthwhile; on commit to a pickup the
agent announces a claim to its teammate (Section~\ref{sec:coord}).

\paragraph{Plans and failure handling.} An intention is achieved by the first
applicable plan in the library, with the rest as fallback (e.g.\ the PDDL
\texttt{go\_to} plan falls back to BFS path following). Movement tolerates a few
failed moves (dynamic obstacles) before declaring a tile soft-blocked so the
next BFS routes around it. Shutdown is close-safe: tearing the agent down
mid-plan (e.g.\ between benchmark runs) cannot crash the process.

\paragraph{Partial observability and limits.} The agent therefore has full
knowledge of the static layout but only partial, decaying knowledge of parcels
and opponents; when no parcel is worth pursuing it explores the
least-recently-seen spawner.

\paragraph{Crate sensing and deterministic push.} Crates (pushable obstacles
on the Sokoban-style maps) are now modelled. They are sensed dynamically from
\texttt{sensing.crates}, stored in the belief base and mirrored onto the graph
as dynamic occupancy, so every BFS detours around them automatically (the map
parser strips a trailing \texttt{!} from tile types and marks type~\texttt{5}
as a pushable zone). Pushing has no dedicated API: when a worthwhile target (a
free parcel, or while carrying a delivery tile) is unreachable \emph{because of}
a sensed crate, the option generator emits a single \texttt{push\_crate} option
--- but only if the push is legal (the tile behind the crate is an empty
type-5 zone, the approach tile is reachable) and a one-tile simulation shows the
target becomes reachable. The deterministic \texttt{PushCrate} plan then walks
to the approach tile and issues exactly one \texttt{move} into the crate; a
rejected push aborts cleanly (one attempt, never a loop) so the normal plans and
BFS take over. The feature is on by default and inert when no crate is sensed.
